\documentclass{article}
\title{Pratice Problem 2.9}
\author{CSAPP}
\usepackage{booktabs}
\usepackage{amsmath}
\usepackage{parskip}

\begin{document}
    \maketitle

    \begin{center}
        \begin{tabular}{cccc}
        \texttt{R}  &   \texttt{G}  &   \texttt{B}  &   \texttt{Color}\\
        \midrule
        \texttt{0}  &   \texttt{0}  &   \texttt{0}  &   \texttt{Black}\\
        \texttt{0}  &   \texttt{0}  &   \texttt{1}  &   \texttt{Blue}\\
        \texttt{0}  &   \texttt{1}  &   \texttt{0}  &   \texttt{Green}\\
        \texttt{0}  &   \texttt{1}  &   \texttt{1}  &   \texttt{Cyan}\\
        \texttt{1}  &   \texttt{0}  &   \texttt{0}  &   \texttt{Red}\\
        \texttt{1}  &   \texttt{0}  &   \texttt{1}  &   \texttt{Magenta}\\
        \texttt{1}  &   \texttt{1}  &   \texttt{0}  &   \texttt{Yellow}\\
        \texttt{1}  &   \texttt{1}  &   \texttt{1}  &   \texttt{White}\\
        \end{tabular}
    \end{center}

    \section*{\small Solution A}
    \begin{flushleft}
        Colors are complemented by complementing the values of 
        \texttt{R, G and B}.\newline \texttt{White} is a complement of \texttt{Black}, 
        \texttt{Yellow} is a complement of \texttt{Blue} and \texttt{Magenta} is the complement of \texttt{Green}
        and \texttt{Cyan} is the complement of \texttt{Red}.
    \end{flushleft}  

    \section*{\small Solution B}
    \texttt{Blue \texttt{(001)} | Green \texttt{(010)}} = \texttt{Cyan \texttt{(011)}}
    \newline\texttt{Yellow \texttt{(110)} \& Cyan \texttt{(011)}} = \texttt{Green \texttt{(010)}}
    \newline\texttt{Red \texttt{(100)} $\land$ Magenta \texttt{(101)}} = \texttt{Blue\texttt{(001)}}
\end{document}